\documentclass[11pt, aspectratio=169]{beamer}
\usepackage[style=authoryear-icomp, backend=biber]{biblatex}
\usepackage{lmodern} % Corrige el escalado de fuentes y símbolos
\usetheme{jcar}

%\includeonlyframes{a} % Para compilar solo en la que se está trabajando

%\addbibresource{refs.bib}

\begin{document}
\author[Juan C. Arboleda Rivera]{Juan Camilo Arboleda Rivera}
\title{Poblaciones y muestras}
%\subtitle{}
%\titlegraphic{\includegraphics[width=2cm]{udea}}
\institute[U. de A.]{Instituto de Biología\\
  Facultad de Ciencias Exactas y Naturales\\
  Universidad de Antioquia\\
  juan.arboleda2@udea.edu.co}
\date{v: \today}
%\subject{}

\begin{frame}[plain]
  \maketitle
\end{frame}

\begin{frame}
 \frametitle{Contenidos}
 \tableofcontents
\end{frame}

\section{Población y muestra}

\begin{frame}
%[label=a]
{}{}
 \centering \Large
 Normalmente cuando tomamos datos, tenemos una muestra, no la población
 completa
 \\~\\
 \includegraphics[height=.65\textheight]{población-muestra.png}
\end{frame}

\begin{frame}
%[label=a]
{Diferencia entre Población y Muestra}{}
 \centering
 \includegraphics[height=.7\textheight]{Población_y_muestra.png}
\end{frame}

\begin{frame}
%[label=a]
{Diferencia entre Población y Muestra}{}
 \centering
 Al tomar una muestra aleatoria hacemos una \textbf{estimación} de un valor
 que queremos conocer, pero esta estimación estará influenciada por el azar.
 \includegraphics[height=.75\textheight]{población-vs-muestra}
\end{frame}

\begin{frame}
%[label=a]
{Parámetro vs. Estadístico}{}
 \centering \Large

 \begin{block}{Parámetro}
  Medida numérica que describe una característica de una población completa
 \end{block}

 \vspace{1em}

 \begin{block}{Estadístico}
  Medida calculada a partir de los datos de una muestra
 \end{block}
% \includegraphics[height=.85\textheight]{}
\end{frame}

\begin{frame}
%[label=a]
{Estimador}{}
 \centering \Large

 Un \textbf{estadístico} utilizado para aproximar un \textbf{parámetro} de
 población se denomina \textbf{estimador} del parámetro
% \includegraphics[height=.85\textheight]{}
\end{frame}

\begin{frame}
%[label=a]
{}{}
 \centering
 \includegraphics[width=\linewidth]{estadístico-parámetro.jpeg}
\end{frame}

\begin{frame}
%[label=a]
{}{}
 \centering \Large
 Dos de los valores más importantes en estadística descriptiva son la
 \textbf{media o promedio} y la \textbf{desviación estándar}.\\~\\
 \includegraphics[height=.75\textheight]{Desviación_estándar.png}
\end{frame}

\begin{frame}
%[label=a]
{}{}
 \centering \Large
 La distinción entre Población y Muestra es muy importante en
 estadística y Ciencia de datos\\~\\

 Cuando se toma una \textbf{muestra} de una \textbf{población} y se calculan
 la media de la muestra ($ \bar{x} $) y la desviación estándar de la muestra
 ($ s $), estos valores se consideran \textbf{estimadores} de la media
 poblacional ($ \mu $) y la desviación estándar poblacional ($ \sigma $)
% \includegraphics[height=.85\textheight]{}
\end{frame}

\begin{frame}
%[label=a]
{}{}
 \centering
 \includegraphics[height=\textheight]{población-muestra2.png}
\end{frame}

\begin{frame}
%[label=a]
{Media o Promedio}{}
 \centering
 La \textbf{media muestral} se representa mediante $ \bar{x} $

 \[ \scalebox{1.5}{$ \displaystyle
 \bar{x} = \dfrac{\sum_{i=1}^{n} x_{i} }{n}
 $}
 \]

 La \textbf{media poblacional} se representa mediante $ \mu $
 {
 \[ \scalebox{1.5}{$ \displaystyle
 \mu = \dfrac{\sum_{i=1}^{N} x_{i} }{N}
 $}
 \]
 }
% \includegraphics[height=.85\textheight]{}
\end{frame}

\begin{frame}
%[label=a]
{Desviación estándar}{}
 \centering
 \textbf{Desviación estándar muestral ($ s $):} Se usa cuando los datos
 provienen de una muestra representativa de la población.\\
 \vspace{-1em}

 \[ s = \sqrt{\dfrac{\sum_{i=1}^{n} (x_i - \bar{x})^2}{n - 1}} \]

 Donde $ x_{i} $ es cada valor, $ \bar{x} $ es la media de la muestra y $ n
 $ es el número total de observaciones.

 \pause

 \textbf{Desviación estándar poblacional ($ \sigma $):} Se usa cuando se
 cuenta con los datos de toda la población.\\
 \vspace{-1em}

 \[ \sigma = \sqrt{\dfrac{\sum_{i=1}^{N}(x_i - \mu)^2}{N}} \]

 Donde $ x_{i} $ es cada valor, $ \mu $ es la media de la población y $ N $
 es el tamaño total de la población.
% \includegraphics[height=.85\textheight]{}
\end{frame}

\begin{frame}
%[label=a]
{Ejercicio}{}
 \centering \Large
 Calcular la desviación estándar de las siguientes variables en Python\\~\\
 \begin{tabular}{|c|c|}
 \hline
 x & y \\
 \hline
 7 & 2 \\
 \hline
 5 & 20 \\
 \hline
 5 & 15 \\
 \hline
 6 & 7 \\
 \hline
 8 & 9 \\
 \hline
 \end{tabular}
% \includegraphics[height=.85\textheight]{}
\end{frame}

\section{Estadística inferencial y Pruebas de hipótesis}

\begin{frame}
%[label=a]
{¿Qué es la estadística inferencial?}{}
 \centering \Large

 La estadística inferencial es una rama de la estadística que se enfoca en
 hacer conclusiones y generalizaciones sobre una población a partir de la
 información obtenida de una muestra de la misma. Para hacer esta
 inferencia, se aplican técnicas estadísticas como el cálculo de intervalos
 de confianza o la realización de pruebas de hipótesis.
% \includegraphics[height=.85\textheight]{}
\end{frame}

\begin{frame}
%[label=a]
{Pruebas de hipótesis}{}
 \centering \Large
 Una prueba de hipótesis es un método estadístico que sirve para decidir si
 una afirmación sobre una población es real o se debe al azar, basándose en
 los datos de una muestra
% \includegraphics[height=.85\textheight]{}
\end{frame}

\section{El valor p}

\begin{frame}
%[label=a]
{Pruebas de hipótesis}{}
 \centering \Large
 El fundamento de la mayoría de pruebas estadísticas usadas en ciencia está
 en responder a la pregunta:\\~\\

 Dado un supuesto (llamado \textbf{Hipótesis nula})\\
 ¿Qué tan probable es obtener este resultado por azar?\\~\\

 \pause

 Esta probabilidad se denomina \textbf{valor p}
% \includegraphics[height=.85\textheight]{}
\end{frame}

\begin{frame}
%[label=a]
{}{}
 \centering
 El \textbf{valor p} o \textbf{p-value} se puede interpretar como la
 probabilidad de obtener un resultado por azar\\
 \vspace{1em}
 \includegraphics[width=\linewidth]{turtle-p-value.jpeg}
 \blfootnote{Bispo, G. S. C. et al. Unveiling the secrets of leatherback
 turtle (\textit{Dermochelys coriacea}) mating system: polygamy, polygyny,
 and sex ratio dynamics of a Brazilian rookery. Biological Journal of the
 Linnean Society 145, blaf026 (2025).
 }
\end{frame}

\begin{frame}
%[label=a]
{Ejemplo: La migración de la polilla Bogong (\textit{Agrotis infusa})}{}
 \centering
 \includegraphics[width=\linewidth]{polillas-bogong.png}
 \blfootnote{Dreyer, D. et al. Bogong moths use a stellar compass for
 long-distance navigation at night. Nature 643, 994–1000 (2025).
 }
\end{frame}

\begin{frame}
%[label=a]
{\centering Test de Rayleigh}{}
 \centering

 \includegraphics[height=.75\textheight]{polar-histogram.png}\\

 \begin{description}
 \item[Hipótesis nula:] Los datos están distribuidos de manera completamente
  uniforme alrededor de un círculo (no hay una dirección preferida)
 \end{description}

\end{frame}

\begin{frame}
%[label=a]
{}{}
 \centering
 \includegraphics[height=\textheight]{migración-bogong.png}
\end{frame}

\begin{frame}
%[label=a]
{Ejemplo}{}
 \centering
 Los colibrís Picoespada (\textit{Ensifera ensifera}) miden (sin contar su
 pico) $ 13.5 \pm 0.3 $
 cm$ ^{*} $\\

 \vspace{.5em}
 \begin{columns}
  \column{.45\linewidth}
   \includegraphics[height=.65\textheight]{Sword-billed_Hummingbird.jpg}
  \column{.55\linewidth}
   \includegraphics[height=.65\textheight]{distribution-ensifera-ensifera.png}
 \end{columns}

 \vfill
 \scriptsize
 $ ^{*} $La desviación estándar no ha sido reportada en ningún
  artículo, fue inferida a partir del rango de tamaños reportado que es
  entre
  13 y 14 cm.
\end{frame}

\begin{frame}
%[label=a]
{}{}
 \centering

 Supongamos que esta especie tiene 20000 individuos\\~\\
 Haremos simulaciones computacionales para ver cómo se comporta un muestreo
 de esta población
% \includegraphics[height=.85\textheight]{}
\end{frame}

%\appendix
%\AtNextBibliography{\footnotesize}
%\begin{frame}[plain, allowframebreaks]{Bibliografía}{}
% \bibitemsep 1em
% \printbibliography[heading=none]
%\end{frame}
\end{document}
